%% kapa: exam=1 round=23 year=2012 period=2 form=A subject=appraisal_law subject_ko=부동산관계법규 seq=7 number=47
%% source: Q-Net 감정평가사 기출문제 (cst00309) artlSeq=5052580 file=raw/1008/5052580/2012년 제23회 감정평가사 제1차시험 기출문제(2교시 A형).pdf
%% conversion: kordoc 4.24.51 -> tools/gen1.py (auto); review=proofread-sonnet; flags=none
\begin{question}{47}
국토의 계획 및 이용에 관한 법령상 용도지구에 관한 설명으로 옳은 것은?
\begin{choices}
  \item 용도지구란 토지의 이용 및 건축물의 용도 등에 대한 용도구역의 제한을 완화하여 적용함으로써 용도구역의 기능을 증진시키기 위하여 도시기본계획으로 결정하는 지역을 말한다.
  \item 복합개발진흥지구는 주거기능, 공업기능, 유통ㆍ물류기능 및 관광ㆍ휴양기능 외의 기능을 중심으로 특정한 목적을 위하여 개발ㆍ정비할 필요가 있는 지구를 말한다.
  \item 집단취락지구는 녹지지역ㆍ관리지역ㆍ농림지역 또는 자연환경보전지역 안의 취락을 정비하기 위하여 필요한 지구를 말한다.
  \item 시ㆍ도지사는 지역여건상 필요하면 법에서 정하고 있는 용도지역 또는 용도구역의 행위제한을 완화하는 용도지구를 신설하여 도시관리계획으로 결정할 수 있다.
  \item 일반미관지구는 중심지미관지구 및 역사문화미관지구 외의 지역으로서 미관을 유지ㆍ관리하기 위하여 필요한 지구를 말한다.
\end{choices}
\end{question}
\answer{5}
